\documentclass[12pt]{article}
\usepackage[margin=1in]{geometry}

% disable page number
\pagestyle{empty}

\begin{document}


My interest in Biomedical Engineering started with my first ever trip to the radiology department of a hospital. I was 16 at the time, and my doctor had me undergo an MRI scan to confirm a diagnosis of rhinitis. Despite how nervous I felt in the noisy, narrow tube for the MRI, I was shocked by the fine slices of my brain that the MRI reconstructed. Ever since then, I've had a deep appreciation for imaging technology, and I want to create new imaging modalities and processing techniques. Because of this, I chose Biomedical Engineering during my undergraduate study at Boston University (BU). My research and coursework confirmed my passion for biomedical imaging, and I would like to pursue research in biophotonics and its applications at Georgia Tech \& Emory University's joint Biomedical Engineering Ph.D. program.
\\

My first research experience in the Dennis Lab, led by Dr. Allison Dennis at BU, was a mixture of wet lab and image processing, which gave me a good foundation for future imaging research. I joined the lab in my Sophomore year and developed methods for wide-field in vivo multiplexed imaging using NIR quantum dots (QD). My main task was to design a fluorophore demixing pipeline and quantify the concentration of each fluorophore. To accomplish this, I adapted a multivariate curve resolution algorithm on MATLAB and validated it with a PDMS phantom setup. I then worked with my graduate student mentor to see if this method would work in a mouse model. To prepare for in vivo studies, I performed lipid-PEG functionalization and encapsulation to transfer our QDs into water. These QDs were injected to image lymph nodes, and my algorithm was tested to demix images. This multifaceted project led to publishable results: I first presented the preliminary results of my work in a poster session at BMES 2019 conference, and my mentor and I are preparing a manuscript for submission. Using this pipeline, we can improve biomolecular specificity and sensitivity over conventional fluorescence imaging.
\\

Overall, my time at the Dennis Lab had a great impact on me. Working on my own has taught me how to independently survey the literature, and the cross-disciplinary work has taught me to be agile in skill acquisition and flexible in troubleshooting. More importantly, this experience has taught me to be resilient in the face of research challenges. Meanwhile, though I had the chance to implement computational methods in imaging at the Dennis Lab, I felt I needed to explore this topic more rigorously. During my Junior year, I took a graduate-level course with Dr. Lei Tian named Computational Optical Imaging. For the class project, I worked on resolving 3D locations of emitters from a single 2D image using double helix and airy beam point-spread functions. This eye-opening class showed me a new design paradigm that integrates optical hardware and computational algorithms. I found it so fascinating to use this design approach to extend new imaging capabilities—such as de-motion blur, super-resolution, and extended depth of field—beyond traditional imaging methods.
\\

From this class, I realized that I needed to enhance my background in signal/image processing to have a better foundation in this subject, leading me to declare a minor in Electrical Engineering. I believe that a solid understanding of the mathematics behind modern processing techniques gives me a deeper perspective on my field of interest. While taking classes to fulfill my degree, I have kept exploring more bioimaging applications through research. This year I joined the BOAS Lab, led by Dr. David Boas at BU, to study the vasculature of post-mortem human brains with neurodegenerative diseases using optical coherence tomography. In this project, I have contributed to the validation and optimization of the automatic vessel segmentation using a Frangi filter. I have gained coding experience implementing algorithms in a shared computing cluster and exposed myself to various neuroimaging techniques. 
\\

The experiences I gained from both research and classes have developed my interest in designing new imaging and processing techniques for biological systems. While exploring this realm requires interdisciplinary expertise in biomedical engineering, Georgia Tech \& Emory's exemplary faculties and its commitment to innovation make it the ideal environment for me to conduct my graduate study. I would be able to learn a great deal if given the opportunity to work with Dr. Shu Jia. His recent work on scaling down the light field microscopy to image the single cell in 3D with high resolution and low computation cost is intriguing to me, and I believe that his novel design scheme will lead to more functional biology studies. I am looking forward to working in this area and study techniques that drive the forward and inverse model design. I would also like to follow Dr. Stanislav Emelianov's work on photoacoustic imaging coupled with nanoparticles, specifically in NIR-II imaging and PA-guided therapy. I believe my prior experience working with QDs can lead me forward in these topics. Last but not least, I would be very interested in working with Dr. Francisco Robles. His research on applying UV microscopy to label-free hematology and hyperspectral imaging is important in translating the design into real clinical use and low-resources settings, which I am hoping to work on. While I have mentioned several faculty members whose work is particularly interesting to me, I anticipate getting a closer look at several other research teams in which to potentially pursue Ph.D. research.
\\

In summary, my goals as a graduate student are to research novel imaging instrumentation and processing pipelines using computational methods. I also want to explore the application of these techniques in the biological system. In the long run, I would like to work as a research scientist in industry or a research institute to keep pushing the limits of imaging technology. I believe that the skills I have gained from my previous experiences have prepared me well for graduate study and research. I look forward to starting a fruitful career at Atlanta and having the opportunity to contribute to the field of biomedical imaging.
\\


\end{document}
